% Created 2024-12-30 Mon 10:51
% Intended LaTeX compiler: pdflatex
\documentclass[11pt]{article}
\usepackage[utf8]{inputenc}
\usepackage[T1]{fontenc}
\usepackage{graphicx}
\usepackage{longtable}
\usepackage{wrapfig}
\usepackage{rotating}
\usepackage[normalem]{ulem}
\usepackage{amsmath}
\usepackage{amssymb}
\usepackage{capt-of}
\usepackage{hyperref}
\author{Hans Halvorson}
\date{January 9, 2025}
\title{What is this thing, \emph{determinism}?}
\hypersetup{
 pdfauthor={Hans Halvorson},
 pdftitle={What is this thing, \emph{determinism}?},
 pdfkeywords={},
 pdfsubject={},
 pdfcreator={Emacs 28.2 (Org mode 9.5.5)}, 
 pdflang={English}}
\begin{document}

\maketitle
\tableofcontents


\section{Introduction}
\label{sec:org58a1aa1}
\begin{itemize}
\item A word of thanks to Klaas Landsman
\end{itemize}

\section{Introduction}
\label{sec:org0e474f4}
\begin{itemize}
\item Brief overview of deterministic theories.
\item Key questions to explore:
\begin{itemize}
\item What are the defining features of a deterministic theory?
\item How does determinism relate to physical laws and initial conditions?
\item Are there examples of controversial or borderline deterministic theories?
\end{itemize}
\end{itemize}

\section{Determinism: attractive and repulsive}
\label{sec:orge2fcec0}
\begin{itemize}
\item Philosophers have been repulsed by determinism because of its
consequences for human action and moral responsibility
\item Physicists have been attracted by determinism because of \ldots{}
\end{itemize}

\section{Philosophy against natural philosophy}
\label{sec:org974db07}
\begin{itemize}
\item Compatibilism might actually be true. But I suspect that it is
also motivated by anti-empiricism.
\item Γνῶθι σεαυτόν
\end{itemize}



\section{Determinism: from world to theory, and back again}
\label{sec:org1756b0a}
\begin{itemize}
\item Carnap
\item Lewis
\end{itemize}

\section{Possible worlds analyses}
\label{sec:org2fcf17d}
\begin{itemize}
\item Earman, \uline{Primer of Determinism}
\end{itemize}



\section{Melia's distinction}
\label{sec:org12b27ef}


\section{Definition}
\label{sec:orgb587d8e}

\textbf{Definition:} We say that a theory \(T\) is \textbf{Landsman-deterministic}
  just in case for any two models \(W\) and \(W'\) of \(T\), and initial
  segment \(U\subseteq W\), if \(f:W\to W'\) and \(g:W\to W'\) are
  isomorphisms such that \(f|_U=g|_U\), then \(f=g\).


\section{What is a theory?}
\label{sec:org1b19d31}



\section{Ideas and Concepts}
\label{sec:org1faebd6}
\begin{itemize}
\item Theoretical Definitions
\begin{itemize}
\item Notes on definitions of determinism in literature.
\end{itemize}
\item Connections to Physical Theories
\begin{itemize}
\item Examples: Newtonian mechanics, quantum mechanics, general relativity.
\end{itemize}
\end{itemize}

\section{PDF Annotations}
\label{sec:org925e247}
\begin{itemize}
\item Below are links to PDFs for reference. Use org-noter to annotate them.
\end{itemize}

** Key References
\begin{itemize}
\item \href{file:///Users/hhalvors/bibliotek/Earman, John/primer-determinism.pdf}{nil}
\item \href{file:///Users/hhalvors/bibliotek/Belot, Gordon/Belot-NewWorkCounterpart-1995.pdf}{nil}
\item \href{file:///Users/hhalvors/bibliotek/Belot, Gordon/Determinism and ontology.pdf}{nil}
\item \href{file:///Users/hhalvors/bibliotek/Melia, Joseph/Melia-HolesHaecceitismTwo-1999.pdf}{nil}
\item Halvorson, Closing the hole argument
\item Landsman, Re-opening the hole argument
\item \href{file:///Users/hhalvors/bibliotek/Montague, Richard/deterministic-theories.pdf}{nil}
\item Recent stuff by Pooley, Read, etc.?
\item Menon and Read, Some remarks on recent formalist responses to
the hole argument \href{file:///Users/hhalvors/bibliotek/Read, James/remarks-hole.pdf}{nil}
\item \href{file:///Users/hhalvors/bibliotek/Landsman, Klaas/bohmian-deterministic.pdf}{nil}
\item \url{https://philpapers.org/browse/determinism}
\item \href{file:///Users/hhalvors/bibliotek/Cudek, Franciszek/15BB8552-BC46-11EE-8419-F108F5211B22.pdf}{nil}
\item \href{file:///Users/hhalvors/bibliotek/Brighouse, Carolyn/determinism-modality.pdf}{nil}
\item \href{file:///Users/hhalvors/bibliotek/Brighouse, Carolyn/understanding-indeterminism.pdf}{nil}
\end{itemize}


** Additional Resources
\begin{itemize}
\item \href{file:///Users/hhalvors/bibliotek/additional\_reading\_1.pdf}{Additional Reading 1}
\begin{itemize}
\item Notes:
\end{itemize}
\item \href{file:///Users/hhalvors/bibliotek/additional\_reading\_2.pdf}{Additional Reading 2}
\begin{itemize}
\item Notes:
\end{itemize}
\end{itemize}

\section{Questions for Discussion}
\label{sec:orgf460118}
\begin{itemize}
\item What is the role of initial conditions in deterministic theories?
\item Can a theory be deterministic at one level and indeterministic at another?
\item How does determinism connect with predictability?
\end{itemize}

\section{To-Do}
\label{sec:org60e0268}
\begin{itemize}
\item[{$\square$}] Read and annotate PDFs using org-noter.
\item[{$\square$}] Summarize main points from annotations.
\item[{$\square$}] Draft a section on examples of deterministic theories.
\end{itemize}

\section{Notes}
\label{sec:org9760241}
\begin{itemize}
\item Add any additional ideas, thoughts, or observations here.
\end{itemize}
\end{document}